\begin{problem}[Dual numbers and invisible first-order motion]
Describe every $k$-algebra map
\[
\varphi:k[t]\to k[\varepsilon]/(\varepsilon^2).
\]
Show that it has the form $t\mapsto a+b\varepsilon$.  Prove that all maps with fixed $a$
induce the same map on underlying spectra, although they are distinct when $b$ differs.
\end{problem}
\paragraph{Why this problem matters.}
It gives a concrete reason nilpotents are useful: they encode first-order variation that
ordinary points cannot see.
\paragraph{Useful ideas before starting.}
A polynomial-ring map is determined by the image of its variable.
\paragraph{Guided hints.}
The dual numbers have exactly one prime ideal $(\varepsilon)$.  Contract it to $k[t]$.
\begin{solution}
Every element of the dual numbers is uniquely $a+b\varepsilon$, so the image of $t$ can be
chosen arbitrarily in this form and determines a unique $k$-algebra map.  The unique point
of the target spectrum is $(\bar\varepsilon)$.  Its contraction consists of polynomials
$f(t)$ with
\[
f(a+b\varepsilon)\in(\varepsilon).
\]
Modulo $\varepsilon$ this condition is $f(a)=0$, so the contraction is $(t-a)$.  Thus $b$
does not affect the underlying point map, even though different $b$ give different ring maps.
\end{solution}
\paragraph{What happened mathematically.}
The topological map records position $a$; the nilpotent coefficient $b$ records a first-order
direction.
\paragraph{Extensions.}
VI/32 will identify $b$ with a tangent-vector parameter.
\paragraph{Provenance.}
New bridge problem anticipating tangent-space geometry without moving the VI/32 theory here.
